% Section 1: introduction, in the order required by the MSSP Guide for Authors: problem statement; state
% of the art by method family; motivation; principle of the approach; contributions; importance and
% impact. The novelty statement (second paragraph of the contributions part) is fixed wording (2026-10-15).
\section{Introduction}\label{sec:intro}

\paragraph{Problem statement.} Each journal of a railway wheelset carries the axle load into the
axlebox through the rollers of the bearing, the bearing housing, the adapter and, on many bogies, an
elastomeric pad. We call the vertical stiffness of this load path the effective support stiffness
$k_B$. It falls when the bearing fits loosen, when the housing, adapter or pad degrades, or when the
bearing clearance grows. The problem addressed here is to estimate $k_B$, and changes of it between
monitoring periods, from the vertical axle-box acceleration (ABA) of a wagon in normal service, when the
rail roughness that excites the response is not measured.

\paragraph{State of the art.} Four families of methods bear on this problem.
\emph{Bearing fault diagnostics} detect localised rolling-contact damage from the impacts it produces,
using envelope and spectral-kurtosis analysis of high-frequency vibration~\citep{randall2011}, or from
the bearing temperature. They are mature and need no model of the vehicle, but they respond to impacts
at roller-passing frequencies, not to a gradual loss of stiffness, so loose fits or a degraded pad
remain invisible until secondary damage occurs.
\emph{Model-based processing of ABA for the track} uses a vehicle--track model to turn ABA into rail
roughness or track stiffness. Roughness is estimated by inverse filtering or source identification
with a known vehicle transfer function~\citep{pieringer2022,carrigan2023}; track stiffness is estimated
from the P2 resonance of single passes~\citep{carrigan2023} or by matching a library of simulations,
where the irregularity level enters as one parameter~\citep{shen2023}. These methods use the known
vehicle as the instrument; they assume the vehicle parameters, and the targets lie below the
axle-box support band. Per-pass estimates also scatter widely when the roughness spectrum has features
near the resonance used~\citep{carrigan2023}.
\emph{Joint identification of vehicle parameters with unknown excitation} has been developed for road
vehicles, where the excitation is separated from the vehicle by the time lag between axles in a Kalman
filter with unknown input~\citep{liu2024}, by a Bayesian formulation in the frequency
domain~\citep{an2024}, or by modelling the road as a Gaussian random field~\citep{guo2026}. Railway
suspension parameters have been identified as a statistical inverse problem with a surrogate model,
using measured track irregularities as input~\citep{lebel2019,lebelhal}. These methods need either the
measured excitation or a model of its structure (time lag between axles, random field), and they have
been demonstrated for suspension parameters at low frequencies, well below the axle-box support band.
\emph{Output-only identification with an unknown excitation spectrum}, such as operational modal
analysis with a Bayesian FFT likelihood~\citep{au2011} or drive-by monitoring that removes road
roughness by subtracting responses recorded at different speeds~\citep{jin2021}, treats the excitation as a nuisance
in a frequency-domain likelihood. It is statistically efficient, but it assumes a smooth excitation
spectrum near each mode and, in its standard form, a single record, so it does not separate a vehicle
resonance from a roughness spectrum that is itself unknown and changes from pass to pass.

\paragraph{Closest work and literature search.} Table~\ref{tab:novelty} compares the method with the
closest work on axle-box acceleration~\citep{carrigan2023,shen2023,pieringer2022} and with axlebox
bearing-response models~\citep{li2019axlebox,guan2023}, which predict bearing loads and vibration for
a given excitation but do not identify the support. The search covered 11--16 October 2026 and used a
general web search engine; bibliographic databases (Scopus, Web of Science, Google Scholar) could not be
queried directly from the computing environment, and full texts were read where they were accessible,
otherwise abstracts. Search terms combined ``axle box acceleration'' or ``axlebox bearing'' with rail
roughness, track stiffness, bearing stiffness or support stiffness, identification, unknown excitation and
speed; further searches covered joint identification of vehicle parameters and road or track
irregularity, operational modal analysis with unknown excitation spectra, drive-by monitoring at several
speeds, and the debiased Whittle likelihood. No lower date limit was set. Only English-language sources
were searched; Chinese- and Japanese-language work on axlebox bearings was not covered.

\begin{table}[tbp]\centering\scriptsize
\caption{Comparison with the closest work. Entries for other work are taken from the cited papers as far as
they could be read (n.r.: not reported in the parts read).}\label{tab:novelty}
\resizebox{\textwidth}{!}{\begin{tabular}{p{2.2cm}p{2.0cm}p{2.1cm}p{1.8cm}p{2.1cm}p{1.6cm}p{2.0cm}p{1.9cm}p{1.9cm}}
\toprule
Work & Target parameter & Excitation assumption & Speed pooling & Roughness treatment & Wheel coupling &
Measurements needed & Validation type & Uncertainty calibration \\
\midrule
Pieringer and Kropp~\citep{pieringer2022} & Rail roughness & Contact force recovered from ABA; track dynamics known & No (single pass) & Estimated output, via a Hertzian contact model & n.r. & ABA, known track parameters & Simulation & Sensitivity to rail-pad stiffness; no intervals \\
Carrigan and Talbot~\citep{carrigan2023} & Track stiffness, then rail roughness & Unknown roughness, removed after the stiffness fit & No (per pass, along the track) & Deconvolution after a P2 peak fit & Several wheels on one rail & ABA, speed, wheel angle (comb filter) & Simulation and metro field trials & Per-pass scatter; no intervals \\
Shen et al.~\citep{shen2023} & Rail-pad and ballast stiffness, irregularity level & Irregularity as a single level parameter & No (three speeds estimated separately) & One level parameter & n.r. & ABA & n.r. & n.r. \\
Li et al.~\citep{li2019axlebox} & None (forward model of bearing loads) & Prescribed track irregularity & -- & Prescribed input & Vehicle model & -- & Wheelset test rig & -- \\
Guan et al.~\citep{guan2023} & None (forward model of bearing response) & Prescribed wheel polygon and irregularity & -- & Prescribed input & Flexible wheelset & -- & Simulation & -- \\
This paper & Bearing support stiffness $k_{B0}$ (and $c_B$) and its change & Unknown roughness, spectrum with free sub-band levels & Yes, many passes in one likelihood & Nuisance levels shared by a period, profiled & Neighbouring wheels via an infinite Timoshenko rail & ABA, speed, WIM payload, distance-aligned records & Matched-model simulation, with model-error and nonlinear-contact checks & CRB, Wald and profile intervals, simulation calibration, inflation for simultaneous model errors \\
\bottomrule
\end{tabular}}
\end{table}

\paragraph{Motivation.} The support resonance of a freight wheelset, the axlebox vibrating on its
support, lies at a few hundred hertz, inside the band in which ABA is excited by rail roughness. ABA is
already recorded in service for track monitoring, so the support stiffness could be tracked from data
that are collected anyway, without a dedicated test, provided the records meet requirements on bandwidth,
calibration, speed and record alignment (Section~\ref{sec:requirements}). None of the families above provides
this: fault diagnostics do not see stiffness, the track methods assume the vehicle, and the
joint-identification methods need the excitation or work below the support band. The work is
motivated by condition monitoring of freight wagons on Indian Railways. The vehicle and track values
that a model of an Indian wagon needs are not yet available to us and are pending data from the
Research Designs and Standards Organisation (RDSO). The method is therefore developed and assessed
here for a representative European freight wagon, a four-axle Eaos/Eanos-type wagon on Y25 bogies on
ballasted UIC60 track, for which the parameters are published; values that could not be found are swept
over plausible ranges.

\paragraph{Principle of the approach.} The support resonance is fixed in frequency, while the roughness
excitation is fixed in wavenumber and therefore moves in frequency with the train speed. Passes at
different speeds thus see the same resonance through different parts of the roughness spectrum. We pool
many passes in one frequency-domain likelihood in which the roughness level of each wavenumber sub-band
is a free nuisance parameter shared by all passes, payload enters through a weigh-in-motion measurement,
and the leakage of the finite records is handled by the debiased Whittle likelihood~\citep{sykulski2019}.
The forward model couples the instrumented wheelset to its neighbours through the rail, because the
other wheels of the bogie and wagon load the rail at the same time.

\paragraph{Contributions.}
To our knowledge, this is the first method to identify the vertical stiffness of the axle-box bearing
support from multi-speed axle-box acceleration without measuring rail roughness. The support resonance
and the roughness are separated by pooling passes at different speeds and payloads in a single
debiased Whittle likelihood, in which roughness levels are free nuisance parameters. Joint estimation
of vehicle or track parameters with unknown excitation has been studied for track stiffness
\citep{carrigan2023,shen2023} and for road-vehicle suspensions \citep{liu2024,an2024}. There, roughness
is handled by deconvolution, a single level parameter, or axle time lags, and the targets lie below the
axle-box support band. Specifically, this paper contributes:
\begin{enumerate}
\item a multi-pass debiased Whittle likelihood that identifies the effective support stiffness and
damping of a wheelset bearing from ABA spectra of many passes, with the roughness level of each
sub-band a free nuisance parameter and payload from a weigh-in-motion measurement;
\item a forward model in which the instrumented wheelset is coupled to its neighbours through an
infinite Timoshenko rail with the phase factors between wheels kept, and evidence that ignoring this
coupling biases the support stiffness by \IgnoreCouplingForty\%;
\item an uncertainty budget for a representative freight wagon, separating statistical error
(Cram\'er--Rao bounds, matched-model Monte Carlo bias and coverage), parameter error (model-form
sensitivity, sweeps over the published parameter ranges, simultaneous model errors), model-structure error
(nonlinear Hertzian contact, payload-dependent primary suspension, wheel coupling) and measurement and
operating conditions;
\item a procedure for detecting changes between monitoring periods that re-estimates the track
foundation and contact stiffness in each period, with detection limits (nominal and allowing for
simultaneous model errors), false-alarm rates under track drift, and the bearing clearance growth that
the detection limit corresponds to. Periods must be compared within one payload band (empty or laden):
pooling passes of mixed payload is not supported by the present model, because the primary suspension
changes with payload.
\end{enumerate}

\paragraph{Importance and impact.} A loss of support stiffness changes the load path of the bearing and
the dynamic loads on it, and it can develop without the temperature rise or impacts that existing
bearing monitoring detects. Detecting it from axle-box records that instrumented wagons already collect
would add a monitoring channel without new hardware on the wheelset. The statistical basis of the method
gives detection limits and false-alarm rates that can be stated in advance for a given monitoring
design. All results in this paper are from matched-model simulations, in which data and estimator share
the vehicle--track model apart from the deliberate perturbations studied; they establish the method and
its error budget, not its operational performance, which requires a controlled rig or field test with a
known change of support stiffness. The same likelihood applies to other vehicle--track resonances excited
by unmeasured roughness.

Section~\ref{sec:model} describes the vehicle--track model and Section~\ref{sec:method} the
estimator. Section~\ref{sec:results} verifies and demonstrates the method in simulation,
Section~\ref{sec:discussion} discusses the results and their limitations, and
Section~\ref{sec:conclusions} concludes.
